\documentclass[10pt,a4paper]{amsart}
\usepackage[margin=19mm]{geometry}
\usepackage{amsmath,amssymb,mathtools}
\usepackage[scr=rsfs,cal=euler]{mathalfa}
\usepackage{xfrac}
\usepackage[unicode,hidelinks,bookmarks=false]{hyperref}
\newcommand{\ie}{i.e.\ }
\newcommand{\cf}{cf.\ }
\newcommand{\resp}{respectively\ }
\newcommand{\Subsection}{\subsection}
\providecommand{\nobreakdash}{\nobreak\hbox{-}}
\setlength{\emergencystretch}{2em}
\multlinegap0pt
\usepackage{mathtools}
\usepackage{amssymb}
\usepackage{xcolor}
\newtheorem{thm}{Theorem}
\newcommand{\BL}{\operatorname{\mathbf{BL}}}
\newcommand{\Lie}[1]{\mathfrak{#1}}
\newcommand{\bsDs}{\boldsymbol{\mathrm{d}\sigma}}
\newcommand{\bsG}{\boldsymbol{G}}
\newcommand{\bsU}{\boldsymbol{U}}
\newcommand{\bsV}{\boldsymbol{V}}
\newcommand{\bsg}{\boldsymbol{\Lie{g}}}
\newcommand{\bsp}{\boldsymbol{p}}
\newcommand{\bss}{\boldsymbol{\sigma}}
\newcommand{\labs}{\left\vert}
\newcommand{\loc}{\mathrm{loc}}
\newcommand{\lpar}{\left( }
\newcommand{\norm}[1]{\left\Vert #1 \right\Vert}
\newcommand{\rabs}{\right\vert}
\newcommand{\rpar}{\right) }
\newcommand{\supp}{\operatorname{supp}}
\newcommand{\wrt}[1]{\,\mathrm{d} #1}
\title{The Brascamp--Lieb inequality on compact Lie groups \\ and its extinction on homogeneous Lie groups}
\author{Michael G. Cowling, Ji Li, Chong-Wei Liang}
\date{Source: arXiv:2602.10647v1 (CC BY 4.0)}
\begin{document}
\maketitle

\noindent\emph{Source and licence.} The Brascamp--Lieb inequality on compact Lie groups \\ and its extinction on homogeneous Lie groups. Version arXiv:2602.10647v1 (CC BY 4.0), distributed by arXiv under the Creative Commons Attribution 4.0 International licence, http://creativecommons.org/licenses/by/4.0/, as recorded in the arXiv licence metadata for this exact version and shown on the abstract page https://arxiv.org/abs/2602.10647v1. Full licence notice: \texttt{licenses/arxiv-2602.10647-CC-BY-4.0.txt}.\par\smallskip

\noindent\textit{Editorial scope.} Counted unit: the article's thm thm:G-loc-g-loc (source0.tex lines 761--767). The article's complete original proof of it is delivered with it (source0.tex lines 769--831, one \texttt{proof} environment of 63 lines). No mathematics is written, repaired or re-derived: the statement and the proof are the article's own lines, in the article's own notation, and no retained line is substituted or re-flowed.  No further source-local context is needed: every cross-reference of every retained line resolves inside the two delivered spans.  Nothing was omitted from any retained span, so the package carries no omission record.  No retained line contains a citation, so the wrapper carries no bibliography and no bibliography entry.  No retained line contains a figure environment, an image-inclusion command or drawing code, so no image is delivered and the package declares no asset dependency.  Editorial formatting only, in addition: an AMS \texttt{amsart} preamble that restates the article's own declarations and macros used by the retained lines, namely the environments \texttt{thm} on separate counters, together with the standard packages those lines need.  The retained environments print in this document as Theorem 1 in order of appearance, whereas the article prints the same items with the numbering of its own full document.  The complete native source of the article as distributed in the arXiv e-print for this exact version is bundled unchanged as \texttt{sources/194416/source0.tex}.
\begin{thm}\label{thm:G-loc-g-loc}
Let $G$ be a Lie group with Lie algebra $\Lie{g}$, and assume that the Haar measures on $G$ and all $G_j$ and the Lebesgue measures on $\Lie{g}$ and all $\Lie{g}_j$ are normalised such that the Radon--Nikodym derivatives of the exponential mappings $\exp_G$ and $\exp_{G_j}$ take the value $1$ at $0$.
If $(G,\bss ,\bsp)$ is a canonical Brascamp--Lieb datum, then 
\begin{align*}
\BL^{\loc}(G,\bss ,\bsp) = \BL^{\loc}(\Lie{g},\bsDs ,\bsp).
\end{align*}
\end{thm}

\begin{proof}
We show first that
\[
\BL^{\loc}(\Lie{g},\bsDs ,\bsp)\leq \BL^{\loc}(G,\bss ,\bsp) .
\]

Take small neighbourhoods $U_j$ of $0$ in $\Lie{g}_j$ and $V_j$ of $e$ in $G_j$ such that $\exp_{G_j} \colon U_j \to V_j $ is a diffeomorphism.
The Radon--Nikodym derivative of $\exp_{G_j}$,  denoted by $J_{G_j}$, is continuous on $\Lie{g}_j$ and takes the value $1$ at $0$.
For all functions $f_j \in C_c(\Lie{g}_j)$ such that $\supp(f_j) \subseteq U_j$, define $\tilde{f} \in C_c(G) $ by letting 
\begin{equation}\label{eq:def-approx-fn}
\tilde f_j(x)\coloneqq  
\begin{cases}
f_j (\log_{G_j} (x))  & \text{if $x \in V_j$} \\
0                      & \text{otherwise}.  
\end{cases}
\end{equation}
Clearly all functions $g_j \in C_c(G_j)$ such that $\supp(g_j) \subseteq V_j$ arise as $\tilde f_j$ for some $f_j \in C_c(\Lie{g}_j)$ such that $\supp(f_j) \subseteq U_j$. 
We claim that 
\begin{equation}\label{eq:claim}
\| \tilde{f} \|_{L^p({G_j})} 
\leq \sup_{X \in U} J_{G_j}(X)^{1/p_j} \| f_j \|_{L^{p_j} (\Lie{g}_j)}.
\end{equation}
This is obvious if $p=\infty$, while if $1\leq p<\infty$, then
\begin{align*}
\| \tilde {f} \|^p_{L^p({G_j})}
&=\int_{G_j} \labs  f(\log_{G_j} (x)) \rabs ^p\wrt x
=\int_{\Lie{g}_j}  \labs f(X) \rabs ^p J_{G_j}(X)\wrt X,
\end{align*}
by a change of variable.
Our claim \eqref{eq:claim} follows.

Without loss of generality, suppose  that the Brascamp--Lieb inputs $f_j$ are nonnegative.

Now,  $\bsV \coloneqq  V_1 \times \dots \times V_J$ is a small neighbourhood of $e$ in $\bsG$, whence $V \coloneqq  \bsV \cap G$ is a small neighbourhood of $e$ in $G$; likewise $\bsU$ is a small neighbourhood of $e$ in $\bsg$ and so $U \coloneqq  \bsU \cap \Lie{g}$ is a samll neighbourhood of $e$ in $\Lie{g}$.
If the $U_j$ are sufficiently small, then the $V_j$ are small, so $V$ is small and hence $U$ is small, so the restriction of $\exp_G$ to $U$ is a diffeomorphism, and we may write
\begin{align*}
 \int_{\Lie{g}} \prod_j f_j\lpar \mathrm{d}\sigma_j(X)\rpar  \wrt X 
&\leq \sup_{X \in U} J_{\Lie{g}}(X)^{-1}  
	\int_{\Lie{g}} \prod_j f_j\lpar \mathrm{d}\sigma_j(X)\rpar J_{\Lie{g}}(X) \wrt X .
\end{align*}
Further, the exponential maps are diffeomorphisms from $U_j$ to $V_j$ and from $U$ to $V$, and $\exp \circ \mathrm{d}\sigma_j = \sigma_j \circ \exp_G$, so
\begin{align*} 
 \int_{\Lie{g}} \prod_j f_j\lpar \mathrm{d}\sigma_j(X)\rpar J_{\Lie{g}}(X) \wrt X 
&=  \int_{G} \prod_j f_j\lpar \mathrm{d}\sigma_j\lpar \log _{G}(x)\rpar \rpar \wrt x\\
&= \int_{G} \prod_j f_j(\log _{G_{j}}\lpar \sigma_j(x)\rpar )\wrt x\\
&= \int_{G} \prod_j \tilde f_j(\sigma_j(x))\wrt x .
\end{align*}
By the definition of $\BL(G,\bss ,\bsp; \bsV)$,
\begin{align*}
 \int_{G} \prod_j \tilde f_j(\sigma_j(x)) \wrt x 
&\leq\BL(G,\bss ,\bsp; \bsV) \prod_j \norm{  \tilde f_j }_{L^{p_j}(G_j)}  \\
&=\BL(G,\bss ,\bsp: \bsV) \prod_j \sup_{X \in U} J_{G_j}(X)^{1/p_j} \prod_j \norm{  \tilde f_j }_{L^{p_j}(G_j)} .
\end{align*}
Putting everything together, it follows that
\[
\BL(\Lie{g},\bsDs ,\bsp; \bsU) 
\leq \sup_{X \in U} J_{\Lie{g}}(X)^{-1}  \prod_j \sup_{X \in U} J_{G_j}(X)^{1/p_j} \BL(G,\bss ,\bsp: \bsV) .
\]
We pass to the limit as $\bsU$ and $\bsV$ shrink, and deduce that $\BL^{\loc}(\Lie{g},\bsDs ,\bsp) \leq 
\BL^{\loc}(G,\bss ,\bsp)$.

The converse inequality is proved similarly. 
\end{proof}

\end{document}
